\subsection{Cognitive Architectures: Integrating Knowledge, Memory, and Reasoning}
\label{sec:4.1.1}
The most ambitious unifying account in contemporary computational neuroscience — influential, and contested on the question of whether one mechanism can carry this much — is neither a pipeline of modules nor a swarm of independent agents. It is a single loop \autocite{friston2009freeenergy}: the brain generates a prediction, compares it against incoming signal, and updates on the error, repeated continuously, at every level, on every kind of input. Perception is that loop pointed at the world, the brain's best guess so far; memory is the loop's store of priors, replayed; attention is not a separate spotlight but the weight the loop assigns to which errors matter; working memory and executive function are that same weighting held active toward a goal, not a separate warehouse or overseer. One mechanism, aimed at different targets, produces what look from the outside like distinct faculties. That is the account's claim.

Nothing forces an AI system to put its dividing lines where a cognitive architecture puts its modules or its agents, because the brain does not put them there either. A system engineered on the predictive-processing model draws its cognitive processes from what the mechanism needs, not from the shape of an org chart. Section~\ref{sec:4.1.2} works the same point through the prefrontal cortex, where control turns out to run without a controller.

The positive result for a designer is that if a single predictive loop is doing the work, a commitment is not a component that can be located and removed. It is a standing bias on what the loop predicts and what it treats as an error worth correcting, which is a harder thing to excise and a harder thing to verify. Section~\ref{sec:4.2} takes that trade up at the point where it becomes a claim about what not to build.
